% ED26-P17-PRACT-003. Erol--Hendler--Nau 1996, pp. 69--93.
\chapter{Hierarchical Task Networks and Plan Existence}
\label{sec:p17-hierarchical-task-network-reasoning}

An executable sequence can reach the desired state and still fail to carry out
the requested task. Buying an object does not carry out a task of making it;
remaining at home does not carry out a round trip. Hierarchical task-network
planning specifies admissible ways of accomplishing tasks as well as the
actions that change the world. This chapter develops the function-free model
and restriction results of Erol, Hendler and Nau
\cite{ErolHendlerNau1996HTNComplexity}. Their 1994 conference presentation is
separate and delegates its proofs to a report
\cite{ErolHendlerNau1994HTN}. Probabilistic policies and utility optimization
belong to Chapter~\ref{sec:p17-sequential-decisions}; the present question is
whether a permitted, executable task decomposition exists.

\section{Actions, goals and compound tasks}
\label{sec:p17-htn-language}

\begin{hbdefinition}{def:p17-htn-vocabulary}{The finite function-free planning vocabulary}
Use mutually disjoint sets of variables \(V\), constants \(C\), predicate
symbols \(P\), primitive-task symbols \(F\), compound-task symbols \(T\), and
task labels \(\mathcal N\). The sets \(C,P,F,T\) are finite and are part of
the problem description; variables and fresh labels are available indefinitely.
Every predicate and task symbol has a fixed finite arity in that description.
Terms are variables or constants: there are no function symbols.

A state \(s\) is a set of ground atoms. An atom is true exactly when it belongs
to \(s\); a negative literal is true when its underlying atom does not.
The three task forms are
\[
 \operatorname{do}[f(\vec x)],\qquad
 \operatorname{achieve}[\ell],\qquad
 \operatorname{perform}[t(\vec x)].
\]
The first is primitive, the second is a goal task, and the third is compound.
Goal and compound tasks are nonprimitive. Distinct labels distinguish repeated
occurrences of the same task.
\index{hierarchical task network}\index{planning!primitive task}
\index{planning!compound task}\index{planning!goal task}
\end{hbdefinition}

\begin{hbdefinition}{def:p17-htn-operator}{Operators and executable traces}
Each primitive symbol has one operator schema with parameter variables,
a conjunction of precondition literals, and positive and negative effects.
For a ground action \(a\), let \(\operatorname{add}(a)\) and
\(\operatorname{del}(a)\) be sets of ground \emph{atoms}; the latter are the
atoms underlying its negative effects. Define
\[
 \operatorname{apply}(s,a)=
 (s\setminus\operatorname{del}(a))\cup\operatorname{add}(a)
\]
when all preconditions hold, and leave it undefined otherwise. The displayed
set update specifies the result even if an atom occurs in both effect sets.

A plan is a finite sequence of ground primitive actions
\(\sigma=a_1\cdots a_m\). Its execution from \(s_0\) has states
\(s_i=\operatorname{apply}(s_{i-1},a_i)\). It is executable if every step is
defined. An effect-free dummy action has empty preconditions and both effect
sets empty; it still occupies a position in a literal action sequence.
\end{hbdefinition}

\section{Networks and conditions on intermediate states}
\label{sec:p17-htn-networks}

\begin{hbdefinition}{def:p17-htn-network}{Task networks and Boolean constraints}
A task network is \(d=[(n_1:\alpha_1),\ldots,(n_m:\alpha_m);\phi]\), with
distinct labels and a Boolean formula \(\phi\) over these atomic constraints:
\begin{enumerate}
\item equalities between variables or constants;
\item ordering constraints \(n_i\prec n_j\);
\item before-state constraints \((\ell,n_i)\) and after-state constraints
      \((n_i,\ell)\);
\item protection constraints \((n_i,\ell,n_j)\).
\end{enumerate}
Negation and disjunction apply to constraints as well as conjunction. Label
expressions \(\operatorname{first}[A]\) and \(\operatorname{last}[A]\), for a
nonempty finite label set \(A\), denote its earliest and latest occurrences.
A primitive network contains only primitive tasks.
\index{planning!protection interval}\index{task network!constraints}
\end{hbdefinition}

\begin{hbdefinition}{def:p17-htn-completion}{A ground primitive completion}
Choose a grounding of the network variables and a permutation placing its
labelled occurrences in an action sequence. Let \(\pi(n)\) be the position
of label \(n\), starting at 1, and require an executable trace
\(s_0,\ldots,s_m\). Equality compares constant symbols, ordering compares
positions, and first/last expressions evaluate by minimum/maximum. The state
constraints mean
\[
 \begin{array}{c|l}
 (\ell,n)&s_{\pi(n)-1}\models\ell\\
 (n,\ell)&s_{\pi(n)}\models\ell\\
 (n,\ell,n')&s_e\models\ell\text{ for every }\pi(n)\leq e<\pi(n').
 \end{array}
\]
The interval includes the state after the left action and excludes the state
after the right action. An empty interval satisfies its universal condition.
The remaining Boolean structure is evaluated in the usual way. The set
\(\operatorname{comp}(d,I,\mathcal D)\) contains exactly the executable
ground orderings satisfying \(\phi\). For a nonprimitive network this set is
empty. For an empty primitive network, the empty sequence is a completion
exactly when its well-formed remaining constraints hold.
\end{hbdefinition}

Negating a protection constraint asserts that some state in the indicated
interval violates the literal. It does not assert that the opposite literal
holds throughout the interval. A correct algorithm must retain this difference.
Similarly, checking an atom immediately before a task differs from inserting
a goal task whose accomplishment may change the state.

\begin{hbdefinition}{def:p17-htn-method}{Methods and the preservation of enclosing constraints}
A method \((\alpha,d')\) gives one permitted decomposition of a nonprimitive
task. Rename all its variables and labels fresh, unify its head with the
selected occurrence, and apply the resulting most general substitution to
both networks. Replace the occurrence \(n\) by the method's labelled tasks,
whose label set is \(A\), and conjoin the method constraint with the transformed
enclosing constraint. Transform each occurrence referring to \(n\) as follows:
\[
 \begin{array}{c|c}
 n\prec u&\operatorname{last}[A]\prec u\\
 u\prec n&u\prec\operatorname{first}[A]\\
 (\ell,n)&(\ell,\operatorname{first}[A])\\
 (n,\ell)&(\operatorname{last}[A],\ell)\\
 (n,\ell,u)&(\operatorname{last}[A],\ell,u)\\
 (u,\ell,n)&(u,\ell,\operatorname{first}[A]).
 \end{array}
\]
Inside an existing first/last expression, replace the label by its new labels,
and continue this rule under later decompositions. These replacements occur
inside arbitrary Boolean formulas without changing their connectives.

Method bodies use nonempty labelled networks so these endpoints are defined.
A zero-effect accomplishment is represented by an explicit dummy task, as in
the source's treatment of an already true goal. Every method for
\(\operatorname{achieve}[\ell]\) includes the requirement that \(\ell\)
hold after its final subtask. Its already-true method uses one dummy task with
the before-condition \(\ell\).
\index{task network!decomposition}\index{planning!method}
\end{hbdefinition}

\begin{hbdefinition}{def:p17-htn-solutions}{Domains and finite derivability}
A domain \(\mathcal D=(\operatorname{Op},\operatorname{Me})\) consists of
finite operator and method sets. A problem is \((d,I,\mathcal D)\).
Write \(\operatorname{red}(d,\mathcal D)\) for its possible one-step
decompositions, including the required substitutions. Define
\[
 \begin{split}
 \operatorname{sol}_0(d,I,\mathcal D)&=\operatorname{comp}(d,I,\mathcal D),\\
 \operatorname{sol}_{r+1}(d,I,\mathcal D)&=
 \operatorname{sol}_r(d,I,\mathcal D)\cup
 \bigcup_{d'\in\operatorname{red}(d,\mathcal D)}
       \operatorname{sol}_r(d',I,\mathcal D),\\
 \operatorname{sol}(d,I,\mathcal D)&=\bigcup_{r\geq0}
       \operatorname{sol}_r(d,I,\mathcal D).
 \end{split}
\]
\emph{Plan existence} asks whether this last set is nonempty.
\end{hbdefinition}

\begin{hbtheorem}{thm:p17-htn-finite-derivation}{proved=annex:p17.htn-finite-derivation}{A solution is a finite permitted derivation followed by execution}
A sequence belongs to \(\operatorname{sol}_r(d,I,\mathcal D)\) exactly when
some sequence of at most \(r\) method decompositions reaches a primitive
network of which it is a completion. In particular, membership in
\(\operatorname{sol}\) requires a finite derivation. An infinite sequence
of task expansions is not itself a plan.
\end{hbtheorem}

\section{A complete task derivation and two failures}
\label{sec:p17-htn-example}

\begin{hbexample}{def:p17-htn-preparation}{Preparing and running a protected task}
Use a constant \(w\), unary predicates for an intact guard, readiness and
completion, and corresponding one-variable task/operator schemas. In this
example abbreviate \(g(w),r(w),c(w)\) by \(g,r,c\), and
\(\mathit{setup}(w),\mathit{run}(w),\mathit{disturb}(w),\mathit{shortcut}(w)\)
by their displayed action names. The initial state is \(I=\{g\}\).
The four ground operator instances are:
\[
 \begin{array}{c|c|c|c}
 \text{action}&\text{precondition}&\operatorname{add}&\operatorname{del}\\\hline
 \mathit{setup}&g&\{r\}&\varnothing\\
 \mathit{run}&r&\{c\}&\varnothing\\
 \mathit{disturb}&\top&\varnothing&\{g\}\\
 \mathit{shortcut}&g&\{c\}&\varnothing
 \end{array}
\]
The compound task \(\mathit{job}(w)\) has the following ground instance of
its sole method schema:
\[
 [a:\operatorname{do}[\mathit{setup}],\ b:\operatorname{do}[\mathit{run}];
 a\prec b\ \land\ (a,g,b)\ \land\ (b,c)].
\]
The schema is obtained by replacing \(w\) consistently by its parameter
variable. In the initial network, place
\(j:\operatorname{perform}[\mathit{job}(w)]\)
beside \(z:\operatorname{do}[\mathit{disturb}]\), with no ordering between
the two. The method supplies the permitted order of its own children, not an
order excluding all possible interleavings with \(z\).
\end{hbexample}

\begin{hbtheorem}{thm:p17-htn-example}{proved=annex:p17.htn-example}{Execution, protection and method-relative success}
For the example, \(\mathit{setup},\mathit{run},\mathit{disturb}\) is a
solution. The sequence \(\mathit{setup},\mathit{disturb},\mathit{run}\)
is executable and ends with \(c\) true, but violates the protected interval
and is not a solution. For the job-only initial network, the one-action
sequence \(\mathit{shortcut}\) also reaches \(c\), but cannot be derived
from its permitted method.

A recursive method \(\mathit{loop}\) decomposing only into itself has no
finite solution. Adding a method from \(\mathit{loop}\) to a dummy primitive
gives a solution, although a search that always chooses the recursive method
can continue indefinitely.
\end{hbtheorem}

Failure under a method library is a failure to obtain an authorized derivation
under that library and its constraints. It need not be physical impossibility.
Likewise, the absence of a result from unbounded search is not a certificate
that no plan exists.

\section{Restriction classes and input size}
\label{sec:p17-htn-restrictions}

\begin{hbdefinition}{def:p17-htn-restrictions}{Primitive, propositional, regular and totally ordered problems}
A problem is \emph{primitive} when its initial network is primitive, and
\emph{propositional} when variables are absent. It is \emph{regular} when
each initial or method network has at most one nonprimitive task and that task
is ordered relative to every other task in its network. Primitive tasks may
precede or follow it. A right-linear method, with the sole nonprimitive task
last, is a special case.

A problem is \emph{totally ordered} when every initial and method network
specifies a total order of its tasks. Having totally ordered methods alone leaves the initial
network free to have unordered roots whose descendants interleave.
Bounds on the number of tasks in a network refer to the initial network and
method bodies in the input description, not to the sizes of all expanded
networks. A \emph{fixed-domain} problem varies the initial state, constants
and initial network while keeping the finite operator and method schemas fixed.
\index{HTN planning!regular}\index{HTN planning!totally ordered}
\end{hbdefinition}

Let \(L\) be the input length in the source's explicit syntax encoding:
argument positions, operator lists and constraint formulas are written out.
Let \(b\) be the number of relevant ground atoms: those appearing in the
initial state or in a ground instance of an operator or network constraint.
Other atom coordinates cannot affect execution or any required constraint
and are omitted from the state descriptions in the complexity arguments.
The function-free vocabulary gives \(b\leq2^{\operatorname{poly}(L)}\), and
there are \(2^b\) world states. In propositional problems \(b\) is polynomial
in \(L\). For a fixed domain, the arities in its schemas are fixed and their
ground predicate and task instances are polynomially numerous in the varying
input. An exponential state space and
an unbounded decomposition space are different sources of difficulty.

\begin{hbtheorem}{thm:p17-htn-semidecision}{proved=annex:p17.htn-semidecision}{General plan existence is semidecidable}
There is a procedure that eventually finds a plan for every solvable finite,
function-free HTN problem. The procedure is not required to terminate on
unsolvable problems.
\end{hbtheorem}

\begin{hblemma}{thm:p17-htn-cfl-intersection}{proved=annex:p17.htn-cfl-intersection}{The exact language-theoretic undecidability premise}
Given two context-free grammars, it is undecidable whether their generated
languages have a common word. The result persists for grammars over a common
binary alphabet, in Chomsky normal form, whose languages exclude the empty
word. Here a context-free grammar has finitely many terminals, nonterminals
and productions, each production replacing one nonterminal by a finite word.
\end{hblemma}

\begin{hbtheorem}{thm:p17-htn-undecidable}{proved=annex:p17.htn-undecidable}{Undecidability despite binary ordered methods}
Plan existence is undecidable even for propositional problems with at most
two tasks in the initial network and in every method body, and with every
method body totally ordered. Thus, together with the preceding theorem, the
problem is strictly semidecidable. The initial two compound tasks need not be
ordered relative to each other.
\end{hbtheorem}

\begin{hbtheorem}{thm:p17-htn-fixed-undecidable}{proved=annex:p17.htn-fixed-undecidable}{A fixed domain can still encode undecidable instances}
There is a fixed function-free domain, with variables and only totally ordered
methods of at most two tasks, for which plan existence is strictly
semidecidable. Its varying initial facts can encode two context-free grammars.
\end{hbtheorem}

\begin{hbdefinition}{def:p17-htn-levels}{A decreasing task level}
A \(k\)-level mapping assigns every ground task a number in
\(\{0,\ldots,k\}\), gives every primitive task level zero, and strictly
decreases from a ground method head to each task in its ground body.
\end{hbdefinition}

\begin{hbtheorem}{thm:p17-htn-levels}{proved=annex:p17.htn-levels}{A finite expansion bound from acyclicity}
Plan existence is decidable for a problem possessing such a level mapping.
Every decomposition branch has depth at most \(k\); finite branching then
bounds the entire decomposition tree. A repeated task with no decreasing
level receives no termination guarantee from this result.
\end{hbtheorem}

\section{The source complexity bounds}
\label{sec:p17-htn-complexity}

The following are the 1996 restriction bounds, keeping gaps between upper and
lower bounds where the source keeps them
\cite[Theorems~4--8]{ErolHendlerNau1996HTNComplexity}. Double exponential time
means a bound \(2^{2^{\operatorname{poly}(L)}}\).

\begin{hbdefinition}{def:p17-htn-trace-summary}{A finite summary retaining interval truth}
For a nonempty execution trace \(\rho=s_0a_1s_1\cdots a_ms_m\), let
\(\Lambda(s)\) be the set of relevant ground literals true in \(s\), and define
\[
 \begin{array}{ll}
 A(\rho)=\displaystyle\bigcap_{j=0}^{m}\Lambda(s_j),&
 B(\rho)=\displaystyle\bigcap_{j=1}^{m}\Lambda(s_j),\\
 C(\rho)=\displaystyle\bigcap_{j=0}^{m-1}\Lambda(s_j),&
 D(\rho)=\displaystyle\bigcap_{j=1}^{m-1}\Lambda(s_j).
 \end{array}
\]
An intersection over no states is the set of all relevant ground literals. The summary
\(\tau(\rho)\) consists of the four boundary states
\(s_0,s_1,s_{m-1},s_m\), these four literal sets, and the Boolean flag
\(m=1\). It records whether a literal holds throughout an interval, not only
whether a chosen protected set was requested.
\end{hbdefinition}

\begin{hblemma}{thm:p17-htn-summary}{proved=annex:p17.htn-summary}{Composition and exact checking from finite summaries}
The summary of a concatenation of compatible nonempty traces is computable
from their summaries. Given summaries for a totally ordered list of task
blocks, every source constraint involving their first/last endpoints can be
evaluated exactly, including negated protection constraints. With \(b\)
ground atoms, there are at most \(2^{12b+1}\) candidate summaries.
\end{hblemma}

\begin{hbtheorem}{thm:p17-htn-total-order}{proved=annex:p17.htn-total-order}{Totally ordered HTN bounds}
For totally ordered problems, plan existence is EXPSPACE-hard and belongs to
double exponential time. In the propositional case it is PSPACE-hard and
belongs to EXPTIME. These pairs of bounds are not assertions of matching
completeness at either endpoint.
\end{hbtheorem}

\begin{hbtheorem}{thm:p17-htn-regular}{proved=annex:p17.htn-regular}{Regular problems and fixed domains}
Regular plan existence is EXPSPACE-complete, and regular propositional plan
existence is PSPACE-complete. The respective lower bounds hold even for
totally ordered problems with one compound-task symbol and input networks
and method bodies of at most two tasks.

With a fixed regular domain and variables, plan existence belongs to PSPACE;
some fixed regular domains make it PSPACE-complete.
\end{hbtheorem}

\begin{hbtheorem}{thm:p17-htn-primitive}{proved=annex:p17.htn-primitive}{Scheduling and binding without decomposition}
Primitive plan existence is NP-complete, remains NP-complete with total
ordering alone or with the propositional restriction alone, and is decidable
in polynomial time when both restrictions are imposed.
\end{hbtheorem}

\section{Comparison with state-goal planning}
\label{sec:p17-htn-strips}

\begin{hbdefinition}{def:p17-htn-strips}{The comparison representation}
A function-free STRIPS-style instance has a finite constant/predicate
vocabulary, finite operator schemas, an initial state and a conjunction of
goal literals. Its solutions are all executable action sequences whose final
state satisfies the goal. No additional task-decomposition requirement applies.

When a translation introduces a bookkeeping action \(\mathit{done}\), define
the observation map \(h\) to erase that action and retain every original
ground action. Erasure is stated explicitly and is not called symbol renaming.
\end{hbdefinition}

\begin{hbtheorem}{thm:p17-htn-strips-embedding}{proved=annex:p17.htn-strips-embedding}{A polynomial regular HTN encoding of state-goal planning}
A STRIPS-style instance has a polynomial-size regular, totally ordered HTN
encoding with one compound-task symbol. Its solutions are exactly the original
solutions followed by one effect-free \(\mathit{done}\) action. Under \(h\),
the two sets of observable solutions agree, including an originally empty plan.
\end{hbtheorem}

\begin{hbtheorem}{thm:p17-htn-expressivity}{proved=annex:p17.htn-expressivity}{Existence and solution-language separation}
There is no computable translation of all function-free HTN instances into
finite function-free STRIPS-style instances preserving plan existence in both
directions.

Moreover, some fixed propositional HTN instances have the nonregular solution
language \(\{a^m b^m:m\geq1\}\). Every finite function-free STRIPS-style
instance has a regular solution language. Thus even a noncomputable choice
of STRIPS instance cannot represent every HTN solution language by symbol
renaming; allowing the erasure of designated bookkeeping actions does not
remove this obstruction. Regular HTN \emph{problems} in the source's sense
need not have regular \emph{solution languages}.
\end{hbtheorem}

These claims distinguish an existence-preserving translation from a weaker
condition checking only the images of source plans. The latter does not by
itself rule out extra target plans. They also distinguish the observed action
sequence from bookkeeping used by an encoding. The practical contribution of
an HTN is its admissible decomposition structure, not a promise that every
hierarchical search will be faster or terminate.
